\section{Demonstration Results}
\label{sec:results}

This section reports qualitative results obtained during prototype
demonstration sessions. The reported sessions exercised the
calibrated load-cell subsystem against a known reference mass, the
DHT11 and MQ135 sensors under indoor ambient conditions, the
ESP32-CAM running both HTTP endpoints concurrently, and the Flutter
web build polling the bridge once per second. The reported numbers
correspond to prototype-level demonstrations rather than to
controlled longitudinal deployments, and the limitations of the
current evaluation are documented in Section~\ref{sec:discussion}.

\subsection{Weight Subsystem Performance}

After a calibration factor of 20.626 was derived from a single-purpose
calibration sketch using a 1\,kg reference mass, the integrated
load-cell subsystem was exercised across several known masses inside
the cabinet. With the cabinet stationary on a level surface,
repeated 10\,kg reference loads were measured between 9.8\,kg and
10.0\,kg. Residual variation is attributable to platform flexure
under load and to load-cell creep.

Two qualitative behaviours were confirmed across sessions. First,
product placement and removal events were detected reliably for
items with a stable bearing surface and a mass of approximately
50\,g or greater. Second, very light objects (folded paper sheets,
single coins) were not consistently detected because their signal
fell within the noise band produced by ambient vibration. A
representative reading near the detection threshold was 0.97\,g
(Listing~\ref{lst:bridgelog}), which corresponds to the subsystem
correctly observing a small perturbation that did not match any
registered product.

\subsection{Environmental Monitoring}

During a ten-minute indoor session without active climate control,
the DHT11 reported temperatures between 27.4\,$^\circ$C and
27.8\,$^\circ$C and relative humidity between 61\,\% and 63\,\%. The
MQ135 reported analog values between 41 and 46 over the same window,
with brief excursions toward 46 when a hand was held close to the
sensor. These observations confirm that the firmware responds to
changes in the cabinet's air composition and reinforce the design
position that MQ135 readings should be interpreted as trend
indicators rather than as calibrated gas concentrations.

\subsection{End-to-End Application Behaviour}

Fig.~\ref{fig:app} reproduces the application during a representative
session. Panel (a) shows the Dashboard rendering the four sensor
cards, a registered banana product, and the active alert summary.
Panel (b) shows the Camera screen overlaying the live MJPEG stream
with the output of the pixel classifier, which on the captured frame
reported $\sigma = 30.6\,\%$ labelled \emph{Spotting --- light
spotting, still consumable}. Panels (c)--(e) show the Products screen
with the registered banana and a three-day-until-expiration band, the
Alerts screen aggregating the two active alerts, and the Settings
screen with the bridge URL and Firebase status.

Listing~\ref{lst:bridgelog} reproduces a short window of bridge log
output. Each \texttt{[serial]} line corresponds to a JSON record
emitted by the Arduino and translated into the bridge's internal
state. During the same window the mobile application polled the
bridge once per second and received a representative response of
\begin{lstlisting}
{"temperature":27.7,"humidity":63.0,
 "gasValue":42,"weight":0,
 "updatedAt":1781692341,"online":true}
\end{lstlisting}
which the Dashboard rendered without intervening cache, confirming
the full data path from the microcontroller through the bridge to
the mobile dashboard.

\subsection{Multi-QR Recognition}

The tiled multi-QR routine was exercised on printed A4 sheets of QR
stickers placed at varied positions and orientations within the
ESP32-CAM's field of view. Up to four distinct QR codes were
recovered from a single capture under cabinet illumination, with
per-frame runtime under 200\,ms on a midrange Android device. In
practice the routine eliminates the need for sequential single-item
scans during product registration, addressing the workflow
penalty that motivated this work (Section~\ref{sec:algo-qr}).

\subsection{Visual Ripeness Estimation}

The pixel-level classifier produced consistent spot-coverage values
across captures of the same banana taken within seconds of each
other. As bananas progressed from a uniformly yellow surface toward
heavier spotting over a multi-day informal observation window, the
reported percentage increased monotonically and the reported stage
label progressed from \emph{Fresh} to \emph{Spotting} and toward
\emph{Spoiling}. Because the classifier maintains no state across
frames, the same input image always produces the same output, which
supports auditability and provides a comparison point against
machine-learning-based detectors.
